% !PRINTED 712
\weberpage{712}
%<<CONT
\weq{2}{g_1 h_1 = \mathop{\mathrm{Lim}}\limits_{s=1} (s-1)\,\Phi_1(s),}
\weq{3}{\Phi_1(s) = \frac{1}{1-q^{-s}}\;\prod^{p_1}\frac{1}{(1-p_1^{-1/2\,sf})^{e}}\;\prod^{p_2}\frac{1}{(1-p_2^{-sf})^{1/2\,e}},}
worin sich die beiden Productzeichen auf alle Primzahlen $p_1$ und
$p_2$ beziehen. Mit Benutzung der oben erklärten Zeichen $f_1$,
$e_1$ kann man dafür auch setzen:
\weq{4}{\Phi_1(s) = \frac{1}{1-q^{-s}}\;\prod^{p}\frac{1}{(1-p^{-sf_1})^{e_1}},}
und das Productzeichen auf \gesperrt{alle} Primzahlen $p$ erstrecken.

Die Bedeutung von $g_1$ ergiebt sich aus \S.\,190, (11).

Nun wollen wir aus der Gruppe $C$ der Charaktere $\chi$ einen
Theiler $C_1$ aussondern, der aus allen den Charakteren $\chi_1$
bestehen soll, für die
\weq{5}{\chi_1(-1) = +1}
ist. Zunächst ist klar, dass diese Charaktere $\chi_1$ eine Gruppe
$C_1$ bilden. Diese Gruppe ist aber nicht mit $C$ identisch, weil
es gewiss einen Charakter $\chi_0$ giebt, für den $\chi_0(-1) = -1$
ist. Wir brauchen, um einen solchen zu erhalten, nur in \S.\,193,
(6) für $\Theta$ eine primitive $\mu^{\text{te}}$ Einheitswurzel,
oder in \S.\,193, (8) $\varepsilon = -1$ zu setzen. Dann bildet die
Gesammtheit der Charaktere $\chi_0\chi_1$, die alle nicht zu $C_1$
gehören, eine Nebengruppe $C_2$, deren Elemente mit $\chi_2$
bezeichnet werden mögen. Jeder Charakter $\chi$ ist also dann
entweder in $C_1$ oder in $C_2$ enthalten, denn wenn $\chi$ nicht
in $C_1$ vorkommt, so kommt $\chi_0^{-1}\chi$ darin vor, und
folglich $\chi$ in $C_2$. Es ist daher $C_1$ ein Theiler von $C$
vom Index 2.

Die Charaktere $\chi_1$ kann man auf folgende Art bilden:

1) Ist $m$ ungerade, so setzen wir, wie in \S.\,193,
\[ n \equiv c^{\gamma} \pmod{m}, \]
und lassen in
\weq{6}{\chi_1(n) = \Theta^{\gamma}}
$\Theta$ die Gesammtheit der $\tfrac12\mu^{\text{ten}}$
Einheitswurzeln durchlaufen.

2) Ist $m$ gerade, so setzen wir
\[ n \equiv (-1)^{\alpha}5^{\beta}, \]
und erhalten
\weq{7}{\chi_1(n) = \Theta^{\beta},}
worin $\Theta$ wieder die sämmtlichen $\tfrac12\mu^{\text{ten}}$
Einheitswurzeln durchläuft.
